%   % !TEX root = ../../VIII,3_Rahmen-TeX_8-1.tex
%
%
%   Band VIII, 3 N.~??A31
%   Signatur/Tex-Datei: LH_35_09_16_015
%   RK-Nr. 41165
%   Überschrift: [Zu John Mayow, Tractatus de sal-nitro et spiritu nitro-aereo]
%   Modul: Mechanik / Elastizität
%   Datierung: [nach Mitte September 1674]
%   WZ: Keins
%   SZ: Keins
%   Bilddateien (PDF): LH_35_09_16_015_d (insgesamt: eins)
%
%
\selectlanguage{ngerman}%
\frenchspacing%
%
\begin{ledgroupsized}[r]{120mm}%
\footnotesize%
\pstart%
\noindent\textbf{Überlieferung:}%
\pend%
\end{ledgroupsized}%
\begin{ledgroupsized}[r]{114mm}%
\footnotesize%
\pstart%
\parindent -6mm%
\makebox[6mm][l]{\textit{L}}%
Notiz zu \textsc{J. Mayow}, \textit{Tractatus de sal-nitro et spiritu nitro-aereo}, in: \textsc{Ders.}, \textit{Tractatus quinque medico-physici}, Oxford 1674:\cite{01997}
LH~XXXV~9,~16 Bl.~15.
Ein Zettel (7 x 2,5 cm).
Eineinhalb Seiten;
Textfolge: Bl.~15~v\textsuperscript{o}, 15~r\textsuperscript{o}.%
\pend%
\end{ledgroupsized}%
%
\vspace{5mm}%
\begin{ledgroup}%
\footnotesize%
\pstart%
\noindent\textbf{Datierungsgründe:}
In seinem Brief an G.\,C. Schelhammer\protect\index{Namensregister}{\textso{Schelhammer} (Schelhammerus), Günther Christoph 1649\textendash1716} von Mitte September 1674 vermerkte Leibniz, er warte gerade auf ein englisches Buch \textit{de respiratione foetus in utero} (\textit{LSB} III,~5 N.~I, S.~3.19\textendash4.1\cite{01293}).
Damit war die gleichnamige Abhandlung gemeint,\cite{01296} die J.~Mayow\protect\index{Namensregister}{\textso{Mayow}, John 1643\textendash1679} 1674 in seinen \textit{Tractatus quinque medico-physici} veröffentlicht hatte und Leibniz vermutlich aus einer Besprechung in den \textit{Philosophical Trans\-ac\-ti\-ons} (20. Juli 1674, S.~101\textendash113\cite{01295}) kannte.
Zum Sammelband gehörte auch Mayows \textit{Tractatus de sal-nitro et spiritu nitro-aereo}\cite{01997}, an den die vorliegende Notiz N.~3 anknüpft.
Unter der Annahme, dass Leibniz das erwartete Buch auch empfing, kann er ab Mitte September 1674 diese Abhandlung gelesen und N.~3 verfasst haben.
\pend%
\pstart
Auf Mayows Sammelband wurde Leibniz ferner von F.~Schrader\protect\index{Namensregister}{\textso{Schrader}, Friedrich 1657\textendash1704} in einem Brief vom 28. November (8. Dezember) 1681 aufmerksam gemacht (vgl. \textit{LSB} III,~3 N.~299, S.~522.5\textendash9; 523.16\textendash18\cite{01294}).
Dies könnte ihn erneut dazu veranlasst haben, sich mit Mayows \textit{Tractatus de sal-nitro et spiritu nitro-aereo} zu beschäftigen.
Es ist ohnehin wahrscheinlich, dass N.~3 auf die frühen Achtziger Jahre zurückgeht, als Leibniz sich intensiv mit der Theorie der Elastizität und Festigkeit befasste, wie zahlreiche Texte im vorliegenden Band bezeugen. 
\pend%
\end{ledgroup}%
%
\selectlanguage{latin}%
\frenchspacing%
%
%
%\newpage%
 \vspace{8mm}
%\pstart 
%\noindent
%\pend
\count\Bfootins=1200
\count\Afootins=1200
\count\Cfootins=1200
%\pstart
%\centering
%\pend
%\count\Bfootins=1000
%\count\Cfootins=1000
%
\pstart%
% \vspace*{0.5em}% PR: nur provisorisch
\noindent
%
\lbrack15~v\textsuperscript{o}\rbrack\ %	%	Blatt 15v
%
\edtext{}{%
{\xxref{LH35_09_16_015v_1}{LH35_09_16_015v_2}}%
{\lemma{Corpus \lbrack...\rbrack\ flexiles}\Cfootnote{%
Siehe J.~\textsc{Mayow}, \cite{01997}\textit{Tractatus de sal-nitro et spiritu nitro-aereo}, S.~73\textendash77; in: \textsc{Ders.}, \textit{Tractatus quinque medico-physici}, Oxford 1674.}}}%
Corpus\edlabel{LH35_09_16_015v_1}
rigidum variis modis flectitur,\protect\index{Sachverzeichnis}{corpus rigidum}\protect\index{Sachverzeichnis}{corpus flexum}
ut vel superficies convexa maneat aeque longa,
concava mutetur;\protect\index{Sachverzeichnis}{superficies convexa}\protect\index{Sachverzeichnis}{superficies concava}
vel ut contra;\protect\index{Sachverzeichnis}{flexio corporis}
vel ut ambae maneant aeque longae ut ante
(\protect\vphantom)%
curvedine\protect\index{Sachverzeichnis}{curvedo} tantum
\edtext{\lbrack mutata\rbrack}{%
\lemma{mutata\protect\vphantom()}\Bfootnote{%
\textit{L~ändert Hrsg.}}}
%
et incluso\lbrack\protect\vphantom()\rbrack,
%
vel ut ambae mutentur. % Dieser Fall nicht bei Mayow, p. 73 sqq.
Mayow\protect\index{Namensregister}{\textso{Mayow}, John 1643\textendash1679}
\textit{de Nitro} c.~6.%
%	%	%	%	ACTHUNG GETRIXT: Folgende Cfootnote bezieht sich auf die Abbildung Fig.1
\edtext{}{\lemma{\lbrack\textit{Fig.~1}\rbrack}\killnumber\Cfootnote{%
Siehe a.a.O., die Abbildung \textit{Fig.~8} auf der Tafel \textit{Tab.~1} am Bandende.}}
%
\pend
 \vspace{1.5em}%
  \centerline{\includegraphics[width=0.35\textwidth]{gesamttex/edit_VIII,3/images/LH_35_09_16_015_d.pdf}}%
  \vspace*{0.0em}%
  \centerline{\lbrack\textit{Fig.~1}\rbrack}% \hspace*{16mm}
  \label{LH_35_09_16_015v_Fig.1}%
\newpage
%
%
\pstart%
\noindent
\textit{a} \edtext{\lbrack chorda\protect\index{Sachverzeichnis}{chorda tensa}\rbrack}{%
\lemma{corda}\Bfootnote{\textit{L~ändert Hrsg.}}}
\quad%
\textit{b} baculus\protect\index{Sachverzeichnis}{baculus}
\quad%
\textit{ab} baculi
\quad
Chorda vel firmiter utrinque alligata, vel ab una parte soluta\lbrack,\rbrack\
baculi flexiles.\edlabel{LH35_09_16_015v_2}
[15~r\textsuperscript{o}]
\pend%
%
%
\pstart%
Multa notanda habet Mayow\protect\index{Namensregister}{\textso{Mayow}, John 1643\textendash1679}
de varietate Elasticorum solidorum\protect\index{Sachverzeichnis}{solidum elasticum} inter tendendum.%
\protect\index{Sachverzeichnis}{tensio corporis}
\pend%
\count\Bfootins=1200
\count\Afootins=1200
\count\Cfootins=1200
%
% PR: Ende des Stückes auf Blatt 15r